
在本章中，我们将研究 $\R^{2n}$ 的这样一类\trSymplecticDiffeomorphism{}所%
构成的群 $\mathcal{D}$：它们由具有\trCompactSupport{}的\trHamiltonianFunction{}%
之\trTimeDependent\trHamiltonianVectorField{}所生成。特别地，我们将循着%
H.~\trHofer{} [116] 的工作，在 $\mathcal{D}$ 上构造一种出人意料的%
\textbf{\trBiInvariantMetric{}}（bi-invariant metric）$d$。借助生成%
\trHamiltonianFunction{}的\trOscillation{}来定义元素 $\psi\in\mathcal{D}$ 的%
\trEnergy{} $E(\psi)$ 之后，\trMetric{} $d$ 便由 $d(\varphi,\psi)=E(\varphi^{-1}\psi)$%
来定义。它是 $C^{0}$ 型的。验证性质「$d(\varphi,\psi)=0$ 当且仅当%
$\varphi=\psi$」并非易事，其依据是\trAction{}原理。一方面，\trMetric{} $d$%
与第 3 章中%
引入的\trCapacity{}函数 $c_{0}$ 密切相关，从而也与诸\trPeriodicOrbit{}密切%
相关；另一方面，它又与定义在 $\R^{2n}$ 的诸子集上且满足 $e(U)\geq c_{0}(U)$%
的一个特殊\trCapacity{} $e$ 密切相关。这个\trCapacity{} $e$ 称为%
\textbf{\trDisplacementEnergy{}}（displacement energy）：$e(U)$ 度量的是%
\trIdentityMap{}与把 $U$ 从其自身挪开的诸 $\psi\in\mathcal{D}$ 所成之集之间的%
\trDistance，此处%
「挪开」意指 $\psi(U)\cap U=\varnothing$。在我们的考察中，元素%
$\psi\in\mathcal{D}$ 之诸\trFixedPoint{}的\textbf{\trActionSpectrum{}}%
（action spectrum）$\sigma(\psi)$ 将扮演一个关键的角色。事实表明，它是 $\R$%
的一个\trCompact{}且\trNowhereDense{}的子集。与第 3 章中对一个固定的%
\trHamiltonianFunction{}所用的简单\trVariational{}技巧相反，我们将设计一个%
\trMinimax{}原理，它同时适用于生成 $\mathcal{D}$ 诸元素的一切%
\trHamiltonianFunction。它挑出一个特定的\trAction{} $\gamma(\varphi)\in\sigma(\varphi)$。%
映射 $\gamma:(\mathcal{D},d)\to\R$ 是 $\mathcal{D}$ 之上%
\textbf{\trActionSpectrum{}\trBundle{}}（action spectrum bundle）的一个连续%
\trSection{}；它是本章的主要技术工具。
它的性质还将用来对 $\mathcal{D}$ 的%
某些元素建立无穷多个\trPeriodicOrbit，并用来描述 $(\mathcal{D},d)$ 的%
\trGeodesic。我们要指出，群 $\mathcal{D}$ 关于\trMetric{} $d$ 的\textbf{\trCompletion}%
（completion）迄今尚不了解。
